\rok{Децембар 2025.}{F}

Први практични тест из Алгоритама и структура података

\zadatak{1}{Quick Sort}
Написати \textit{quickSort} методу Quick Sort алгоритма за сортирање целобројног низа дужине $n$.

\textbf{Додатне напомене за решавање задатка:}
\begin{itemize}
	\item Улаз је несортиран низ целих бројева.
	\item Излаз је сортирани низ целих бројева.
	\item У template-у се налази класа \texttt{RandomArrayGenerator} коју треба искористити за генерисање низа случајних бројева.
	\item У оквиру класе \texttt{QuickSort} написати имплементацију за методу:
	\begin{Verbatim}[fontsize=\small]
private static void quickSort(int[] arr, int lower, int upper)
	\end{Verbatim}
	\item Обезбедити код у Main методи за тестирање алгоритма.
\end{itemize}

\zadatak{2}{Проналажење најдуже подлисте парних бројева}
Написати алгоритам који проналази најдужу подлисту парних елемената. Вредности из пронађене подлисте исписати у конзоли, или да таква подлиста не постоји.

\textbf{Додатне напомене за решавање задатка:}
\begin{itemize}
	\item Захтеван потпис методе: \texttt{public static void findLargestEvenNumbersSubList()}.
	\item Методу писати у оквиру класе \texttt{LinkedList}.
\end{itemize}

\textbf{Потпис методе:}
\begin{Verbatim}[fontsize=\small]
public static void findLargestEvenNumbersSubList()
\end{Verbatim}

\textbf{Пример:}
\begin{itemize}
	\item \textbf{Улаз:} \texttt{LinkedList: 1 -> 0 -> 6 -> 8 -> 4} \\
	      \textbf{Излаз:} \texttt{0 6 8 4}
	\item \textbf{Улаз:} \texttt{LinkedList: 1 -> 5 -> 3} \\
	      \textbf{Излаз:} \texttt{Не постоји подлиста парних бројева}
\end{itemize}

\zadatak{3}{Hot Potato Game -- Odd and Even Problem}
\textbf{Проблем:} У класичној игри "Врући кромпир", играчи у кругу пребацују кромпир док музика свира. Када музика стане, играч који држи кромпир испада. У нашој напредној варијанти, број пролазака кромпира се динамички мења током игре.

\textbf{Правила игре:}
\begin{enumerate}
	\item Има $n$ играча нумерисаних од 1 до $n$.
	\item Почетни број пролазака је $K$.
	\item Правило елиминације: У сваком кругу, кромпир се проследи $K$ пута. Играч код кога се кромпир нађе након $K$-тог проласка испада.
	\item Динамичка промена $K$: Након што играч испадне, нови број пролазака ($K_{novi}$) за следећи круг се рачуна на основу ID-а елиминисаног играча:
	\begin{itemize}
		\item Ако је ID испалог играча паран број: $K_{novi} = K_{stari} + 2$
		\item Ако је ID испалог играча непаран број: $K_{novi} = K_{stari} - 1$
	\end{itemize}
	\item Ограничења за $K$:
	\begin{itemize}
		\item Ако $K$ постане мањи од 1, поставља се на вредност почетног $K$ (из параметра методе).
		\item Пре него што почне нови круг, ако је $K$ већи од броја преосталих играча, узима се остатак при дељењу:
		$$
		K = (K \% \text{broj\_preostalih}) + 1
		$$
	\end{itemize}
	\item Игра се наставља док не остане само један играч -- победник.
\end{enumerate}

\textbf{Додатни захтеви за решавање задатка:}
\begin{enumerate}
	\item Задатак решавати у оквиру методе:
	\begin{Verbatim}[fontsize=\small]
public int PlayHotPotato(int totalPlayers, int initialPasses)
	\end{Verbatim}
	\item Морате користити Queue за представљање круга играча.
	\item Имплементирати динамичко рачунање пролазака по датој формули.
	\item Резултат извршавања алгоритма треба да буде редни број победника у игри.
\end{enumerate}

\textbf{Пример:}
\begin{itemize}
	\item \textbf{Улаз 1:} \texttt{totalPlayers = 4, initialPasses = 2} \\
	      \textbf{Излаз 1:} \texttt{expectedWinner = 3}
	\item \textbf{Улаз 2:} \texttt{totalPlayers = 6, initialPasses = 1} \\
	      \textbf{Излаз 2:} \texttt{expectedWinner = 6}
\end{itemize}

\sablon{Шаблон првог теста децембар 2025. група F}{2025/Decembar/GrupaF/AISP2025_Test1_GrupaF.linq}
